% Front-door note N2 (periods/motives audience) — Phase 2a of the
% corpus accessibility plan. NOT SENT; sends are Phase 4, PI-gated.
% Every identity below is symbolically proved and numerically
% verified to >= 40 digits; provenance flagged inline.
\documentclass[11pt]{article}
\usepackage[T1]{fontenc}
\usepackage[utf8]{inputenc}
\usepackage{amsmath, amssymb, amsthm}
\usepackage[margin=1.1in]{geometry}
\usepackage{hyperref}

\newtheorem{identity}{Identity}
\newtheorem{question}{Question}

\newcommand{\C}{\mathbb{C}}
\newcommand{\Q}{\mathbb{Q}}
\newcommand{\Z}{\mathbb{Z}}

\title{Three exact identities from Dirac spectral zetas on the round
$S^3$,\\ and a question about their motivic homes}
\author{J.~Loutey\thanks{Independent researcher,
\texttt{jloutey@gmail.com}. Developed in an AI-assisted workflow
(implementation and drafting with Anthropic's Claude);\ each identity
below is proved by elementary manipulations reproduced here in full
and verified numerically to at least $40$ digits;\ errors are the
author's. Full write-ups and verification scripts:\
\texttt{https://github.com/jloutey-hash/geovac}.}}
\date{June 2026}

\begin{document}
\maketitle

\noindent Each identity is checkable in roughly ten minutes from the
definitions given;\ the question at the end is the reason for
writing.

\paragraph{Setup.} The Dirac operator on the round unit $S^3$ has
spectrum $\pm(n + \tfrac32)$, $n \ge 0$, with multiplicity
$(n+1)(n+2)$ per sign. Write $m = n + \tfrac32$ and define the
(absolute) spectral Dirichlet series and its parity sectors
\begin{equation*}
  D(s) = \sum_{n \ge 0} 2(n+1)(n+2)\, m^{-s}, \qquad
  D_{\mathrm{even}}(s) = \!\!\sum_{n \ \mathrm{even}}\!\!(\cdots),
  \qquad
  D_{\mathrm{odd}}(s) = \!\!\sum_{n \ \mathrm{odd}}\!\!(\cdots).
\end{equation*}
Let $\zeta(s, a)$ be Hurwitz zeta, $\beta(s) = L(s, \chi_{-4}) =
4^{-s}\bigl(\zeta(s, \tfrac14) - \zeta(s, \tfrac34)\bigr)$ Dirichlet
beta, $G = \beta(2)$ Catalan's constant.

\begin{identity}[parity sectors and a $\chi_{-4}$ bridge]
Since $2(n+1)(n+2) = 2(m^2 - \tfrac14)$, the even-$n$ eigenvalues are
$m = 2(2k + \tfrac34)$ and the odd-$n$ ones $m = 2(2k + \tfrac54)$,
giving
\begin{equation*}
  D_{\mathrm{even}}(s) = 2^{-(s-3)}\zeta(s{-}2, \tfrac34)
                       - 2^{-(s+1)}\zeta(s, \tfrac34), \qquad
  D_{\mathrm{odd}}(s) = 2^{-(s-3)}\zeta(s{-}2, \tfrac54)
                       - 2^{-(s+1)}\zeta(s, \tfrac54),
\end{equation*}
and, from $\zeta(s, \tfrac34) - \zeta(s, \tfrac54)
= 4^s(1 - \beta(s))$ (shift recurrence plus the $\beta$ definition),
\begin{equation*}
  D_{\mathrm{even}}(s) - D_{\mathrm{odd}}(s)
  \;=\; 2^{s-1}\bigl(\beta(s) - \beta(s{-}2)\bigr)
  \qquad (s > 3;\ \text{all } s \text{ by continuation}).
\end{equation*}
At $s = 4$:
\begin{equation*}
  D_{\mathrm{even}}(4) = \tfrac{\pi^2}{2} - \tfrac{\pi^4}{24}
  - 4G + 4\beta(4), \qquad
  D_{\mathrm{odd}}(4) = \tfrac{\pi^2}{2} - \tfrac{\pi^4}{24}
  + 4G - 4\beta(4),
\end{equation*}
so the $\beta$-values cancel in the sum,
$D(4) = \pi^2 - \pi^4/12$, and survive only in the parity
difference, $D_{\mathrm{even}}(4) - D_{\mathrm{odd}}(4)
= 8(\beta(4) - G)$.
\end{identity}

\begin{identity}[vanishing of the spectral zeta at non-positive
integers]
With full degeneracy, $\zeta_{D^2}(s) := \sum_{n} 2(n+1)(n+2)\,
m^{-2s} = 2\zeta(2s{-}2, \tfrac32) - \tfrac12 \zeta(2s, \tfrac32)$
satisfies
\begin{equation*}
  \zeta_{D^2}(-k) \;=\; 0 \qquad \text{for every integer } k \ge 0.
\end{equation*}
The mechanism is the Bernoulli identity $B_{2k+1}(\tfrac32) =
(2k{+}1)/4^k$ (from $B_{2k+1}(\tfrac12) = 0$ and the shift formula),
which forces exact pairwise cancellation between the two Hurwitz
pieces at every non-positive integer. Consequence:\ the
Chamseddine--Connes spectral action of $D$ on $S^3_R$ is
\emph{exactly two-term} --- $\Lambda^3$- and $\Lambda^1$-terms only,
all higher Seeley--DeWitt contributions vanish identically, the
remainder is $O(e^{-\Lambda^2 R^2})$.
\end{identity}

\begin{identity}[all-orders closed form for the scalar heat trace]
For the conformally shifted scalar Laplacian on unit $S^3$
(eigenvalues $n^2$ on the shifted spectrum, integer Jacobi-$\vartheta_3$
sum), the heat trace is
\begin{equation*}
  K_\Delta(t) \;=\; \frac{\sqrt{\pi}}{4}\,\frac{e^{t}}{t^{3/2}}
  \;+\; O\!\bigl(e^{-\pi^2/t}\bigr),
\end{equation*}
i.e.\ \emph{every} Seeley--DeWitt coefficient is
$a_k = 2\pi^2/k!$ in the standard normalization --- the full
asymptotic series is $\pi^2 \cdot \Q[t]$, exactly summable.
\end{identity}

\paragraph{The pattern, and the question.} Across a sizable family
of spectral invariants computed on this geometry, the transcendental
content sorts into exactly two rings:
\begin{itemize}
\item \emph{untwisted} spectral sums (heat traces, spectral actions,
full zeta values:\ Identities 2--3, and e.g.\
$D(4) = \pi^2 - \pi^4/12$) land in the pure-Tate ring
$\bigoplus_k \pi^{2k}\,\Q$ --- consistent with the
Fathizadeh--Marcolli classification of Seeley--DeWitt coefficients
as mixed-Tate periods, but in a strictly smaller sub-ring (no
$\zeta(3)$, no MZVs);
\item the \emph{parity-twisted} sector (Identity 1) lands in
mixed-Tate periods of level $4$ ($\beta$-values, Catalan), i.e.\
cyclotomic over $\Z[i, \tfrac12]$ in the sense of Deligne's
classification at $N = 4$ and Glanois' bases --- and the
even/odd \emph{difference} collapses onto the pure $\beta$-line,
which reads like a Galois descent from level $4$ to level~$2$.
\end{itemize}

\begin{question}
Is there a motivic-Galois statement that \emph{predicts} this
sector-to-ring assignment --- untwisted sums to
$\bigoplus_k \pi^{2k}\Q$, parity-twisted sums to level-$4$
cyclotomic mixed Tate, with the parity difference realizing the
descent --- rather than confirming it case by case? Concretely:\ is
the ``parity acts as $\chi_{-4}$, difference $=$ descent''
observation in Identity~1 an instance of a known functoriality in
the cyclotomic mixed-Tate category, and if so, what else does that
functoriality predict for spectral invariants on
$S^3 = \mathrm{SU}(2)$?
\end{question}

\noindent\emph{Provenance:} Identity 1 is proved above (the three
displayed reductions);\ Identity 2's Bernoulli mechanism and
Identity 3's $\vartheta_3$ computation are short and reproduced in
the full write-ups, with all three verified symbolically
(\texttt{sympy}) and numerically to $\ge 40$ digits against direct
spectral sums.

\end{document}
